\section*{الفصل \ref{ch:b1:isostasy-gravity} --- التوازن الأيزوستاتي وشذوذات الجاذبية}

\begin{solution}{exo:b1:isostasy-gravity:1}
(أ) $\qty{-150}{\milli\gal} = -150\times\qty{e-5}{m.s^{-2}} = \qty{-1.5e-3}{m.s^{-2}}$،
و$\num{1.5e-3}/9.81 = \num{1.53e-4}$ من $g$. (ب) $\qty{1}{\milli\gal} =
\qty{e-5}{m.s^{-2}}$، أي \num{1.02e-6} من $g$؛ و$\qty{1}{\micro\gal} =
\qty{e-8}{m.s^{-2}}$، أي \num{1.02e-9} من $g$: جهاز يميّز الميكروغال يقيس الجاذبية بدقة نحو جزء من مليار.
\end{solution}

\begin{solution}{exo:b1:isostasy-gravity:2}
(أ) عند خط الاستواء تنعدم كلتا الجيبيتين: $\gamma = \qty{9.780327}{m.s^{-2}}$. عند
\ang{45}، $\sin^2\varphi = 0.5$ و$\sin^2 2\varphi = 1$:
$\gamma = 9.780327\times(1 + 0.0026512 - 0.0000058) = \qty{9.806200}{m.s^{-2}}$.
عند القطبين، $\sin^2\varphi = 1$ و$\sin 2\varphi = 0$:
$\gamma = 9.780327\times1.0053024 = \qty{9.832186}{m.s^{-2}}$.
(ب) $9.832186 - 9.780327 = \qty{0.051859}{m.s^{-2}}$، أي
\qty{5186}{\milli\gal}. (ج) $3392/5186 = 0.65$: يفسّر الأثر الطارد المركزي نحو ثلثي الفرق؛ أما التفلطح، الذي يقرّب القطبين من المركز، فيفسّر الباقي.
\end{solution}

\begin{solution}{exo:b1:isostasy-gravity:3}
$r = h\rho_c/(\rho_m - \rho_c) = 3\times2800/500 = \qty{16.8}{km}$. وتصبح \omterm{def:b1:layered-earth:layers}{القشرة} أسمك بمقدار $h + r = \qty{19.8}{km}$ (أي $h\rho_m/(\rho_m - \rho_c) =
3\times3300/500$).
\end{solution}

\begin{solution}{exo:b1:isostasy-gravity:4}
(أ) $1/0.3086 = \qty{3.24}{m}$. (ب) $0.01/0.3086 = \qty{0.032}{m}$: يجب أن تُعرف الارتفاعات بدقة نحو \qty{3}{cm}، ولذلك يسجّل المسح ارتفاع كل محطة بتحديد المواقع بالأقمار الصناعية أو بالتسوية.
\end{solution}

\begin{solution}{exo:b1:isostasy-gravity:5}
(أ) $\rho = 2700\times100/(100 + h)$: \qty{2673}{kg.m^{-3}} عند \qty{1}{km}،
\qty{2621}{kg.m^{-3}} عند \qty{3}{km}، و\qty{2571}{kg.m^{-3}} عند \qty{5}{km}.
(ب) $\rho = (2700\times100 - 1030\times4)/96 = \num{265880}/96 =
\qty{2770}{kg.m^{-3}}$: العمود تحت البحر أكثف. (ج) يقتضي نموذج برات تغيّرات جانبية في الكثافة عند عمق ثابت. وتوجد تلك التغيّرات تحت قاع المحيط: \omterm{def:b1:layered-earth:lithosphere}{فالغلاف الصخري} الفتيّ والحار قرب الحيد أخفّ من الغلاف القديم البارد البعيد عنه (\cref{ch:b1:heat-budget}). أما تحت القارات فعلى العكس، إذ يُلاحَظ أن سمك \omterm{def:b1:layered-earth:layers}{القشرة} يتغيّر مع التضاريس.
\end{solution}

\begin{solution}{exo:b1:isostasy-gravity:6}
(أ) المحطة 1: $g - \gamma = 9.807751 - 9.808911 = \qty{-0.001160}{m.s^{-2}} =
\qty{-116.0}{\milli\gal}$؛ \omterm{def:b1:isostasy-gravity:anomalies}{تصحيح الهواء الحر} $0.3086\times350 =
\qty{108.0}{\milli\gal}$، فإن $\Delta g_{\mathrm{FA}} = \qty{-8.0}{\milli\gal}$؛
\omterm{def:b1:isostasy-gravity:anomalies}{تصحيح بوغيه} $0.112\times350 = \qty{39.2}{\milli\gal}$، فإن
$\Delta g_{\mathrm B} = \qty{-47.2}{\milli\gal}$.
المحطة 2: $g - \gamma = \qty{-535.9}{\milli\gal}$؛ \omterm{def:b1:isostasy-gravity:anomalies}{تصحيح الهواء الحر}
$0.3086\times1850 = \qty{570.9}{\milli\gal}$، فإن $\Delta g_{\mathrm{FA}} =
\qty{+35.0}{\milli\gal}$؛ \omterm{def:b1:isostasy-gravity:anomalies}{تصحيح بوغيه} $0.112\times1850 =
\qty{207.2}{\milli\gal}$، فإن $\Delta g_{\mathrm B} = \qty{-172.2}{\milli\gal}$.
(ب) عند المحطة 2 يكون \omterm{def:b1:isostasy-gravity:anomalies}{شذوذ بوغيه} سالبًا بقوة: تحت الجبال يوجد عجز كتلي، أي جذر من \omterm{def:b1:layered-earth:layers}{القشرة} الخفيفة. ويبيّن \omterm{def:b1:isostasy-gravity:anomalies}{شذوذ الهواء الحر}، الصغير مقارنةً \omterm{def:b1:isostasy-gravity:anomalies}{بتصحيح بوغيه}، أن العجز يوازن الجبال تقريبًا: فهي معوَّضة في معظمها (بنحو $172/207$، أي \qty{83}{\percent}، إذا صحّ نموذج الصفيحة للسلسلة كلها).
\end{solution}

\begin{solution}{exo:b1:isostasy-gravity:7}
(أ) كما في برهان القضية، تُسهم طبقة رقيقة على المسافة $z$ بمقدار $G\rho\,2\pi z\,h\int_0^a s\,(s^2 + z^2)^{-3/2}\,\mathrm ds$، مع توقف التكامل عند $s = a$:
$\int_0^a s\,(s^2 + z^2)^{-3/2}\,\mathrm ds = 1/z - 1/\sqrt{z^2 + a^2}$. ومنه
$g = 2\pi G\rho h\,\bigl(1 - z/\sqrt{z^2 + a^2}\bigr)$، وهي تؤول إلى قيمة الصفيحة عندما $a \to \infty$. (ب) $z/\sqrt{z^2 + a^2} = 0.1$ يعطي
$\sqrt{z^2 + a^2} = 10z$، و$a = \sqrt{99}\,z = \qty{9.95}{km}$ من أجل
$z = \qty{1}{km}$. (ج) الصخر الواقع ضمن نحو عشرة أمثال عمقه تحت المحطة يوفّر \qty{90}{\percent} من الجذب: فتصحيح التضاريس يحتاج إلى التضاريس على نصف قطر من بضعة كيلومترات إلى عشرات الكيلومترات حول كل محطة، لا تحتها وحده.
\end{solution}

\begin{solution}{exo:b1:isostasy-gravity:8}
(أ) $r' = d(\rho_c - \rho_w)/(\rho_m - \rho_c) = 4\times1770/500 =
\qty{14.16}{km}$، فيكون سمك \omterm{def:b1:layered-earth:layers}{القشرة} $35 - 4 - 14.16 = \qty{16.84}{km}$.
(ب) يتنبّأ نموذج إيري بكثافة قشرية واحدة \omterm{def:b1:layered-earth:layers}{بقشرة} محيطية أسمك بأكثر من ثلاث مرات مما هي عليه. \omterm{def:b1:layered-earth:layers}{والقشرة} المحيطية ليست \omterm{def:b1:layered-earth:layers}{قشرة} قارية رقيقة: إنها صخر آخر أكثف، والمستويان على الأرض نوعان من \omterm{def:b1:layered-earth:layers}{القشرة}، كما يوحي \omterm{def:b1:isostasy-gravity:hypsometry}{المنحنى الهيبسومتري}.
\end{solution}

\begin{solution}{exo:b1:isostasy-gravity:9}
(أ) الثانية القوسية الواحدة هي $\pi/(180\times3600) = 1/\num{206265}$ من الراديان:
\qty{5.236}{''} $= \qty{2.538e-5}{rad}$ و\qty{15.885}{''} $=
\qty{7.701e-5}{rad}$. (ب) $a_h = g\theta$: $9.79\times\num{2.538e-5} =
\qty{2.485e-4}{m.s^{-2}} = \qty{24.9}{\milli\gal}$ مرصودة، و $9.79\times\num{7.701e-5} = \qty{75.4}{\milli\gal}$ محسوبة من الجبال المرئية. (ج) $1 - 5.236/15.885 = 0.670$: ألغى التعويض نحو ثلثي جذب الجبال المرئية.
\end{solution}

\begin{solution}{exo:b1:isostasy-gravity:10}
(أ) خارجها تجذب الكرة كما تجذب كتلة نقطية في مركزها، بالمقدار $GM/(x^2 + z^2)$ على الخط الواصل إلى المركز؛ والمركبة الشاقولية هي ذلك مضروبًا في جيب التمام $z/\sqrt{x^2 + z^2}$، وهذا يعطي $GMz/(x^2 + z^2)^{3/2}$.
(ب) $M = \tfrac43\pi(\qty{2000}{m})^3\times(\qty{-200}{kg.m^{-3}}) =
\qty{-6.70e12}{kg}$، وفوق المركز $GM/z^2 = \num{6.674e-11}\times
(\num{-6.70e12})/(4000)^2 = \qty{-2.80e-5}{m.s^{-2}} = \qty{-2.80}{\milli\gal}$.
(ج) $(1 + x^2/z^2)^{-3/2} = 1/2$ يعطي $x = z\sqrt{2^{2/3} - 1} = 0.766\,z =
\qty{3.07}{km}$. (د) نصف عرض المنحنى عند نصف قيمته العظمى يعطي العمق: $z = 1.305\,x_{1/2}$. والجسم الأعمق يعطي شذوذًا أعرض وأكثر انبساطًا.
\end{solution}

\begin{solution}{exo:b1:isostasy-gravity:11}
(أ) $V = \tfrac13\pi R^2 H = \tfrac13\pi(\qty{15e3}{m})^2(\qty{3e3}{m}) =
\qty{7.07e11}{m^3}$، و$\Delta m = (2800 - 1030)V = \qty{1.25e15}{kg}$.
(ب) يقع مركز الكتلة على \qty{0.75}{km} فوق القاع، أي على عمق $d = 5 - 0.75 = \qty{4.25}{km}$ تحت سطح البحر. وهناك يكون شذوذ الجهد $T = -G\Delta m/d$، فإن $N = G\Delta m/(\gamma d) =
\num{6.674e-11}\times\num{1.25e15}/(9.81\times4250) = \qty{2.0}{m}$.
(ج) معظم كتلة المخروط تقع في قاعدته الواسعة السفلى، على مسافات من 5 إلى \qty{16}{km} من النقطة فوق القمة: فمتوسط $1/r$ على الكتلة أصغر من $1/\qty{4.25}{km}$، فتبالغ نقطة الكتلة في تقدير الجهد. وبجمع الأقراص الرقيقة نحصل على نحو \qty{1.1}{m}. (د) يضع جذرٌ من \omterm{def:b1:layered-earth:layers}{قشرة} خفيفة تحته عجزًا من الكتلة نفسها على عمق أكبر؛ فالجهدان يتلاشيان تقريبًا، ويصبح الانتفاخ أصغر بكثير، لكنه ليس صفرًا، لأن العجز يقع أبعد عن السطح من الفائض.
\end{solution}

\begin{solution}{exo:b1:isostasy-gravity:12}
(أ) معدّل الرفع هو $-\mathrm dw/\mathrm dt = w_0\,e^{-t/\tau}/\tau = w/\tau$.
(ب) $w = 818\,e^{-10000/4000} = 818\times0.0821 = \qty{67.1}{m}$، ومعدّل الرفع هو $\qty{67.1}{m}/\num{4000}\,\unit{yr} = \qty{16.8}{mm.yr^{-1}}$، أي أكثر من مرة ونصف القيمة المقاسة. (ج) مع $\tau = 3100$ سنة يكون $w =
818\,e^{-3.226} = \qty{32.5}{m}$ والمعدّل \qty{10.5}{mm.yr^{-1}}؛ ومع $\tau = \num{68000}$ سنة يكون $w = 818\,e^{-0.147} = \qty{706}{m}$ والمعدّل \qty{10.4}{mm.yr^{-1}}. وهذان يقتضيان هبوطين متباينين جدًا. والهبوط غير المعوَّض عجز كتلي من صخر \omterm{def:b1:layered-earth:layers}{الوشاح}: فعلى منطقة واسعة يعطي شذوذ هواء حر قدره نحو $-2\pi G\rho_m w$، أي $-0.0419\times3.3\times706 \approx \qty{-98}{\milli\gal}$ للحالة البطيئة، مقابل \qty{-4}{\milli\gal} للحالة السريعة. ويفصل بينهما مسح للجاذبية (أو ارتفاع الشواطئ القديمة). (د) يقيس ثابت الزمن مدى سرعة عودة \omterm{def:b1:layered-earth:layers}{الوشاح} للتدفّق تحت الهبوط: أي لزوجته.
\end{solution}

\begin{solution}{pb:b1:isostasy-gravity:1}
\textbf{1.} $r = h\rho_c/(\rho_m - \rho_c) = 4\times2700/600 = \qty{18}{km}$.

\textbf{2.} سمك \omterm{def:b1:layered-earth:layers}{القشرة} $35 + 4 + 18 = \qty{57}{km}$؛ وقاعدتها، أي \omterm{def:b1:layered-earth:discontinuity}{موهو}، تقع على $35 + 18 = \qty{53}{km}$ تحت مستوى البحر.

\textbf{3.} الهضبة: \qty{57}{km} من \omterm{def:b1:layered-earth:layers}{القشرة} و$60 - 53 = \qty{7}{km}$ من \omterm{def:b1:layered-earth:layers}{الوشاح}،
$\num{57000}\times2700 + \num{7000}\times3300 = \num{1.539e8} + \num{2.31e7} =
\qty{1.770e8}{kg.m^{-2}}$. المنخفض: $\num{35000}\times2700 + \num{25000}\times3300
= \num{9.45e7} + \num{8.25e7} = \qty{1.770e8}{kg.m^{-2}}$. وهما متساويان، كما
يقتضي \cref{prop:b1:isostasy-gravity:equal-mass}.

\textbf{4.} $p = g\times\text{الكتلة لكل وحدة مساحة} = 9.81\times\num{1.770e8} =
\qty{1.74e9}{Pa}$، أي نحو \qty{1.7}{GPa}، تحت العمودين معًا.

\textbf{5.} يجب أن يعوّض الجذر، بكثافته الأصغر، كتلة الهضبة: فكل متر من الجذر يزيل $\rho_m - \rho_c = \qty{600}{kg.m^{-3}}$ من الكتلة، وكل متر من الهضبة يضيف \qty{2700}{kg.m^{-3}}، إذن $2700/600 = 4.5$ متر من الجذر لازمة. ومع \omterm{def:b1:layered-earth:layers}{قشرة} كثافتها \qty{2900}{kg.m^{-3}} تكون النسبة $2900/400 = 7.25$: فكلما صغر التباين، عمق الجذر أكثر.

\textbf{6.} يتعمّق \omterm{def:b1:layered-earth:discontinuity}{موهو}، أي الزيادة القوية في سرعة الموجات الزلزالية، تحت حافة الهضبة، من \qty{35}{km} تحت المنخفض إلى نحو \qty{53}{km} تحت الهضبة.

\textbf{7.} $\rho = 2700\times100/104 = \qty{2596}{kg.m^{-3}}$.

\textbf{8.} $\num{104000}\times2596.2 = \qty{2.700e8}{kg.m^{-2}}$ و
$\num{100000}\times2700 = \qty{2.700e8}{kg.m^{-2}}$.

\textbf{9.} $2700 - 2596 = \qty{104}{kg.m^{-3}}$، أو $104/2700 =
\qty{3.85}{\percent}$.

\textbf{10.} $\Delta\rho/\rho = \alpha\,\Delta T$ يعطي $\Delta T =
0.0385/\num{3e-5} = \qty{1280}{K}$ متوسطًا على عمود \qty{104}{km} كاملًا. وهذا كافٍ لصهر جزء كبير من العمود: فالحرارة وحدها لا تجعل هذه الهضبة خفيفة بما يكفي، ولا يمكن أن تكون آلية برات هي السبب الرئيس هنا.

\textbf{11.} تبقى أعمدة برات بلا جذر: يبقى \omterm{def:b1:layered-earth:discontinuity}{موهو} على نحو \qty{35}{km} تحت الهضبة، كما تحت المنخفض.

\textbf{12.} \qty{52}{km} أقرب إلى \qty{53}{km} عند إيري منها إلى \qty{35}{km} عند برات: فالهضبة معوَّضة في الأساس بجذر من \omterm{def:b1:layered-earth:layers}{القشرة}.

\textbf{13.} $0.0419\times2.7\times4000 = \qty{452.5}{\milli\gal}$.

\textbf{14.} بلا تعويض (\cref{prop:b1:isostasy-gravity:broad}):
$\Delta g_{\mathrm{FA}} \approx \qty{+452.5}{\milli\gal}$،
$\Delta g_{\mathrm B} \approx 0$.

\textbf{15.} مع التعويض: $\Delta g_{\mathrm{FA}} \approx 0$ و
$\Delta g_{\mathrm B} \approx \qty{-452.5}{\milli\gal}$. ويعطي العجز الكتلي للجذر $-0.0419\times0.6\times\num{18000} = \qty{-452.5}{\milli\gal}$ القيمة نفسها، لأن $\Delta\rho\,r = \rho_c h$.

\textbf{16.} $\sin^2\ang{30} = 0.25$ و$\sin^2\ang{60} = 0.75$:
$\gamma = 9.780327\times(1 + 0.0013256 - 0.0000044) = \qty{9.793249}{m.s^{-2}}$.
ثم $g - \gamma = 9.781085 - 9.793249 = \qty{-1216.4}{\milli\gal}$؛ \omterm{def:b1:isostasy-gravity:anomalies}{وتصحيح الهواء الحر}
$0.3086\times4000 = \qty{1234.4}{\milli\gal}$، فإن
$\Delta g_{\mathrm{FA}} = \qty{+18.0}{\milli\gal}$، و
$\Delta g_{\mathrm B} = 18.0 - 452.5 = \qty{-434.5}{\milli\gal}$.

\textbf{17.} $434.5/452.5 = 0.96$: الهضبة معوَّضة بنحو \qty{96}{\percent} (\omterm{def:b1:isostasy-gravity:anomalies}{شذوذ الهواء الحر} الموجب الصغير يمثّل النسبة \qty{4}{\percent} غير المعوَّضة).

\textbf{18.} لا. قرب مركز هضبة واسعة تعمل كل طبقة شاذة كصفيحة، وجذبها لا يتعلّق بعمقها: لا يهمّ إلا العجز الكتلي لكل وحدة مساحة، وهو نفسه، $\rho_c h$، في النموذجين. ويختلفان فقط قرب الحواف، حيث يغيّر عمق الكتل المعوِّضة شكل الشذوذ.

\textbf{19.} بعيدًا عن الحافة على أيّ من الجانبين يكون \omterm{def:b1:isostasy-gravity:anomalies}{شذوذ الهواء الحر} قرب الصفر. وعلى بضعة كيلومترات داخل الحافة تكون الهضبة تحت المحطة قريبة وتمتد بعيدًا على جانب واحد فقط: فهي ما زالت تجذب بقرب قيمتها الكاملة للصفيحة. أما الجذر، على \qty{35}{km} أو أكثر تحتها، فيُرى من هناك كجسم ينتهي تقريبًا مستقيمًا تحتها: إنه يجذب بأقل قليلًا من نصف أثره. والفرق قمة موجبة. وعلى بضعة كيلومترات خارج الحافة تجذب الهضبة قليلًا جدًا، بينما يبقى الجذر يجذب بقرب نصف أثره: منخفض سالب. فالقطاع زوج موجب-سالب عبر كل حافة، كما في شكل الفصل.

\textbf{20.} \omterm{def:b1:layered-earth:layers}{القشرة} الآن $57 - 1 = \qty{56}{km}$ سمكًا، أي $\qty{21}{km}$ أكثر من \omterm{def:b1:layered-earth:layers}{قشرة} المنخفض. وفي التوازن ينقسم هذا الفائض إلى $h' + r'$ مع $r' = 4.5\,h'$: $h' = 21/5.5 = \qty{3.818}{km}$ و$r' = \qty{17.18}{km}$.

\textbf{21.} ترتفع الهضبة الآن إلى \qty{3818}{m}.

\textbf{22.} كانت الصخور التي على السطح الآن على عمق \qty{1}{km} تحت السطح القديم، أي عند \qty{3000}{m}؛ وهي الآن عند \qty{3818}{m}: فقد ارتفعت بمقدار \qty{818}{m}، أي $1000\times\rho_c/\rho_m$.

\textbf{23.} كل كيلومتر يُزال يخفض السطح بمقدار $1 - 2700/3300 =
0.182$ من الكيلومتر، فإنزاله بمقدار \qty{4}{km} يقتضي
$4/0.182 = \qty{22}{km}$ من التعرية، أي السمك الزائد كله $h + r$. ولذلك كانت الصخور التي على سطح سلسلة جبلية قديمة مستوية بالتعرية على عمق نحو عشرين كيلومترًا يومًا ما: فجذور السلسلة العميقة مكشوفة.

\textbf{24.} ينخفض السطح بمقدار $1000 - 818$: لكل كيلومتر من الصخر المُزال، يكون الانخفاض الصافي للسطح
\textbf{\qty[text-series-to-math=true]{182}{m}}.
\end{solution}
